\begin{definition}[Denominator-sensitive adjoint stationary map]\label{mod:parameters-adjoint}
Let $K/F$ be a finite extension of nonarchimedean local fields.  For
$E=F,K$, let
\[
 m_E=2d_E+\varepsilon_E,\qquad
 1<m_E,\qquad \varepsilon_E\leq1,
\]
and put
\[
 C_E=\mathcal O_E/\mathfrak p_E^{d_E},
 \qquad
 V_E=\mathfrak p_E^{d_E+\varepsilon_E}/\mathfrak p_E^{m_E}.
\]
Thus $C_E$ is the Lean type
\texttt{StationaryCoefficientQuotient $E$ $d_E$}.  Let
\[
 P_{K/F}:V_K\longrightarrow V_F
\]
be the additive truncated norm map constructed from a
\texttt{NormPolynomialPrecision} certificate.

Fix exact additive-character data $\psi_E$ and denominators
$\gamma_E\in E^\times$ satisfying
\[
 \operatorname{ord}_E(\gamma_E)=m_E+n_E(\psi_E).
\]
Finite duality gives the additive equivalence
\[
 \Phi_{E,\gamma_E}:C_E\xrightarrow{\ \sim\ }
   \operatorname{AddChar}(V_E,\C),
 \qquad
 \Phi_{E,\gamma_E}(\bar c)(\bar y)
   =\psi_E(cy/\gamma_E).
\]
Both arguments are quotient classes, and evaluation is independent of
representatives modulo $\mathfrak p_E^{d_E}$ and
$\mathfrak p_E^{m_E}$ respectively.

Define the pulled-back character and the denominator-sensitive adjoint by
\[
 P^\vee_{\gamma_F}(\bar c)
   =\Phi_{F,\gamma_F}(\bar c)\circ P_{K/F},
 \qquad
 P^*_{K/F;\gamma_F,\gamma_K}
   =\Phi_{K,\gamma_K}^{-1}\circ P^\vee_{\gamma_F}.
\]
This is an additive homomorphism
\[
 P^*_{K/F;\gamma_F,\gamma_K}:C_F\longrightarrow C_K
\]
and satisfies, for all $\bar c\in C_F$ and $\bar y\in V_K$,
\[
 \Phi_{F,\gamma_F}(\bar c)\bigl(P_{K/F}(\bar y)\bigr)
 =
 \Phi_{K,\gamma_K}
   \bigl(P^*_{K/F;\gamma_F,\gamma_K}(\bar c)\bigr)(\bar y).
\]
Equivalently, for supplied representatives $c,b,y$ such that
$P^*_{K/F;\gamma_F,\gamma_K}([c])=[b]$,
\[
 \psi_F\!\left(
   \frac{c\,\bigl(N_{K/F}(1+y)-1\bigr)}{\gamma_F}\right)
 =
 \psi_K\!\left(\frac{by}{\gamma_K}\right),
\]
where the norm displacement on the left is the certified representative
of the truncated norm class.  Perfectness proves that this adjoint is the
unique additive homomorphism satisfying the displayed pairing identity.

All needed representative invariances are explicit.  If
$c\equiv c'\pmod{\mathfrak p_F^{d_F}}$, then the two adjoint inputs give
the same class.  If their images are represented by $b$ and $b'$, then
\[
 b\equiv b'\pmod{\mathfrak p_K^{d_K}}.
\]
If $y\equiv y'\pmod{\mathfrak p_K^{m_K}}$, both sides of the
representative-level adjoint identity are unchanged.  These are quotient
equalities and character equalities; no field-valued representative of an
adjoint class is chosen.

For two admissible denominators $\gamma,\gamma'$ over the same field, put
\[
 a_{\gamma,\gamma'}=\frac{\gamma'}{\gamma}.
\]
It has order zero, so multiplication by it defines a
representative-independent additive endomorphism
\[
 S_{\gamma,\gamma'}:C_E\longrightarrow C_E,
 \qquad [c]\longmapsto[a_{\gamma,\gamma'}c].
\]
The pairing obeys
\[
 \Phi_{E,\gamma'}\bigl(S_{\gamma,\gamma'}(\bar c)\bigr)
 =\Phi_{E,\gamma}(\bar c).
\]
Consequently adjoints for two denominator pairs satisfy both pointwise and
as additive homomorphisms
\[
 P^*_{K/F;\gamma_F',\gamma_K'}\circ
   S_{\gamma_F,\gamma_F'}
 =
 S_{\gamma_K,\gamma_K'}\circ
   P^*_{K/F;\gamma_F,\gamma_K}.
\]

Suppose now that
\[
 \psi_K=\psi_F\circ\operatorname{Tr}_{K/F}
\]
and that the truncated norm map equals the trace map on $V_K$.  Assume
separately the trace-lattice inclusions at depths
$d_K+\varepsilon_K$ and $m_K$, the quotient-depth bound
\[
 d_K\leq e(K/F)d_F,
\]
and
\[
 \operatorname{ord}_K\!\left(
   \frac{\gamma_K}{\iota(\gamma_F)}\right)=0.
\]
Then the adjoint is exactly the quotient homomorphism
\[
 P^*_{K/F;\gamma_F,\gamma_K}([c])
 =
 \left[
   \frac{\gamma_K}{\iota(\gamma_F)}\,\iota(c)
 \right]
 \in\mathcal O_K/\mathfrak p_K^{d_K}.
\]
This equality is also exported for every supplied integral representative
$c$.  In the common-denominator case
$\gamma_K=\iota(\gamma_F)$ it specializes to
\[
 P^*_{K/F;\gamma_F,\iota(\gamma_F)}([c])=[\iota(c)].
\]
Both specializations assert equality only in the target quotient.

\LeanFile{LanglandsFirstMainLemma/Parameters/Adjoint.lean}
\lean{LanglandsFirstMainLemma.StationaryCoefficientQuotient}
\lean{LanglandsFirstMainLemma.stationaryPairingLeftEquiv}
\lean{LanglandsFirstMainLemma.stationaryPairingLeft}
\lean{LanglandsFirstMainLemma.stationaryPairingLeft_mk_mk}
\lean{LanglandsFirstMainLemma.normPolynomialPullbackCharacter}
\lean{LanglandsFirstMainLemma.normPolynomialAdjoint}
\lean{LanglandsFirstMainLemma.normPolynomialAdjoint_character}
\lean{LanglandsFirstMainLemma.normPolynomialAdjoint_pairing}
\lean{LanglandsFirstMainLemma.normPolynomialAdjoint_unique}
\lean{LanglandsFirstMainLemma.normPolynomialAdjoint_existsUnique}
\lean{LanglandsFirstMainLemma.stationaryPairingLeft_coefficient_representative_independent}
\lean{LanglandsFirstMainLemma.stationaryPairingLeft_variable_representative_independent}
\lean{LanglandsFirstMainLemma.normPolynomialAdjoint_pairing_mk}
\lean{LanglandsFirstMainLemma.normPolynomialAdjoint_input_representative_independent}
\lean{LanglandsFirstMainLemma.normPolynomialAdjoint_representatives_congruent}
\lean{LanglandsFirstMainLemma.normPolynomialAdjoint_variable_representative_independent}
\lean{LanglandsFirstMainLemma.denominatorChangeRatio}
\lean{LanglandsFirstMainLemma.denominatorChangeRatio_ne_zero}
\lean{LanglandsFirstMainLemma.denominatorChangeRatio_ord_eq_zero}
\lean{LanglandsFirstMainLemma.stationaryCoefficientScaledRepresentative}
\lean{LanglandsFirstMainLemma.stationaryCoefficientScale}
\lean{LanglandsFirstMainLemma.stationaryCoefficientScale_mk}
\lean{LanglandsFirstMainLemma.stationaryCoefficientScale_representative_independent}
\lean{LanglandsFirstMainLemma.stationaryPairingLeft_changeDenominator}
\lean{LanglandsFirstMainLemma.normPolynomialAdjoint_changeDenominator}
\lean{LanglandsFirstMainLemma.normPolynomialAdjoint_changeDenominator_hom}
\lean{LanglandsFirstMainLemma.normPolynomialAdjoint_eq_denominatorScaledAlgebraMap}
\lean{LanglandsFirstMainLemma.normPolynomialAdjoint_trace_denominatorRatio_mk}
\lean{LanglandsFirstMainLemma.normPolynomialAdjoint_commonDenominator_mk}
\leanok
\SourceText{Lemma~\ref{lem:truncated-norm-adjoint}.}
\uses{mod:parameters-normpolynomial,mod:lamprecht-finiteduality}
\FormalizationContract The truncated norm has direction
$V_K\to V_F$, while its adjoint has direction $C_F\to C_K$.
The source and target coefficient depths are exactly $d_F$ and $d_K$.
Both denominators, their order proofs, and every denominator-change factor
are explicit inputs.  Existence and uniqueness are derived from the
perfect left pairing rather than stored as fields of an adjoint structure.
All maps are defined on finite lattice quotients, all representative
invariances needed downstream are separate theorems, and the trace and
common-denominator formulas assert quotient equality without choosing
canonical field representatives or strengthening it to equality in $K$.
\end{definition}
